\documentclass[12pt,a4paper]{article}
\usepackage[utf8]{inputenc}
\usepackage[T1]{fontenc}
\usepackage[french]{babel}
\usepackage{geometry}
\usepackage{graphicx}
\usepackage{svg}
\usepackage{xcolor}
\usepackage{hyperref}
\usepackage{fancyhdr}
\usepackage{titlesec}
\usepackage{booktabs}
\usepackage{array}
\usepackage{longtable}
\usepackage{colortbl}
\usepackage{listings}
\usepackage{tcolorbox}
\usepackage{enumitem}
\usepackage{tikz}

% Page geometry
\geometry{margin=2cm}

% Colors
\definecolor{primary}{HTML}{E11D48}
\definecolor{primarydark}{HTML}{BE123C}
\definecolor{secondary}{HTML}{0EA5E9}
\definecolor{success}{HTML}{10B981}
\definecolor{warning}{HTML}{F59E0B}
\definecolor{dark}{HTML}{1E293B}
\definecolor{light}{HTML}{F8FAFC}
\definecolor{gray}{HTML}{64748B}
\definecolor{codebg}{HTML}{1E293B}

% Hyperref setup
\hypersetup{
    colorlinks=true,
    linkcolor=primary,
    urlcolor=secondary,
    pdftitle={Guide DevOps - CNSS},
    pdfauthor={CNSS}
}

% Header/Footer
\pagestyle{fancy}
\fancyhf{}
\fancyhead[L]{\textcolor{gray}{\small Guide DevOps - CNSS}}
\fancyhead[R]{\textcolor{gray}{\small Mars 2026}}
\fancyfoot[C]{\textcolor{gray}{\thepage}}
\renewcommand{\headrulewidth}{0.4pt}
\renewcommand{\footrulewidth}{0pt}

% Section formatting
\titleformat{\section}
    {\Large\bfseries\color{primary}}
    {\thesection.}{0.5em}{}[\color{primary}\titlerule]
    
\titleformat{\subsection}
    {\large\bfseries\color{dark}}
    {\thesubsection}{0.5em}{}

% Code listing style
\lstset{
    basicstyle=\ttfamily\small,
    backgroundcolor=\color{codebg},
    commentstyle=\color{success},
    keywordstyle=\color{secondary},
    stringstyle=\color{warning},
    frame=single,
    rulecolor=\color{gray},
    breaklines=true,
    showstringspaces=false,
    columns=flexible,
    literate={é}{{\'e}}1 {è}{{\`e}}1 {ê}{{\^e}}1 {à}{{\`a}}1 {ù}{{\`u}}1
}

% Custom boxes
\tcbuselibrary{skins,breakable}

\newtcolorbox{infobox}{
    colback=secondary!10,
    colframe=secondary,
    leftrule=4pt,
    rightrule=0pt,
    toprule=0pt,
    bottomrule=0pt,
    arc=0pt,
    breakable
}

\newtcolorbox{successbox}{
    colback=success!10,
    colframe=success,
    leftrule=4pt,
    rightrule=0pt,
    toprule=0pt,
    bottomrule=0pt,
    arc=0pt,
    breakable
}

\newtcolorbox{warningbox}{
    colback=warning!10,
    colframe=warning,
    leftrule=4pt,
    rightrule=0pt,
    toprule=0pt,
    bottomrule=0pt,
    arc=0pt,
    breakable
}

% Badge command
\newcommand{\badge}[2]{%
    \tikz[baseline=(char.base)]{
        \node[fill=#1,text=white,rounded corners=3pt,inner sep=3pt,font=\tiny\bfseries] (char) {#2};
    }%
}

\begin{document}

% Title Page
\begin{titlepage}
    \centering
    \vspace*{2cm}
    
    {\Huge\bfseries\textcolor{primary}{Guide DevOps}\par}
    \vspace{0.5cm}
    {\LARGE Mise en Production\par}
    \vspace{1cm}
    {\Large\textcolor{gray}{Architecture Microservices CNSS}\par}
    
    \vspace{2cm}
    
    \begin{tikzpicture}
        \draw[primary,line width=3pt] (0,0) -- (10,0);
    \end{tikzpicture}
    
    \vspace{2cm}
    
    {\large
    \begin{tabular}{rl}
        \textbf{Application:} & Coopération Technique + Mise en Disponibilité \\[5pt]
        \textbf{Stack:} & Spring Boot · Angular · Python · Docker · Kubernetes \\[5pt]
        \textbf{Version:} & 1.0 \\[5pt]
        \textbf{Date:} & Mars 2026 \\
    \end{tabular}
    }
    
    \vfill
    
    {\Large\bfseries CNSS - Caisse Nationale de Sécurité Sociale\par}
    {\large Tunisie\par}
    
\end{titlepage}

% Table of Contents
\tableofcontents
\newpage

% ============================================
\section{Architecture de l'application CNSS}
% ============================================

L'application CNSS est composée de 15 microservices interconnectés, répartis en plusieurs couches :

\begin{itemize}[leftmargin=*]
    \item \textbf{Frontend} : 2 applications Angular 17
    \item \textbf{Gateway} : NGINX Ingress, API Gateway, Eureka Server
    \item \textbf{Microservices} : 10 services Spring Boot
    \item \textbf{Service IA} : Python FastAPI + Tesseract OCR
    \item \textbf{Infrastructure} : Oracle DB, RabbitMQ, Redis
\end{itemize}

\subsection{Diagramme d'architecture}

\begin{center}
    \fbox{\includesvg[width=0.95\textwidth]{images/architecture-cnss.svg}}
\end{center}

\subsection{Liste des microservices}

\begin{center}
\rowcolors{2}{light}{white}
\begin{longtable}{|c|l|l|c|l|}
\hline
\rowcolor{primary}
\textcolor{white}{\textbf{\#}} & \textcolor{white}{\textbf{Service}} & \textcolor{white}{\textbf{Technologie}} & \textcolor{white}{\textbf{Port}} & \textcolor{white}{\textbf{Image Docker}} \\
\hline
\endfirsthead
\hline
\rowcolor{primary}
\textcolor{white}{\textbf{\#}} & \textcolor{white}{\textbf{Service}} & \textcolor{white}{\textbf{Technologie}} & \textcolor{white}{\textbf{Port}} & \textcolor{white}{\textbf{Image Docker}} \\
\hline
\endhead
1 & eureka-server & Spring Boot & 8761 & ghcr.io/cnss/eureka \\
2 & gateway-service & Spring Cloud Gateway & 8080 & ghcr.io/cnss/gateway \\
3 & auth-service & Spring Boot + JWT & 8089 & ghcr.io/cnss/auth \\
4 & employer-service & Spring Boot & 8081 & ghcr.io/cnss/employer \\
5 & salary-service & Spring Boot & 8082 & ghcr.io/cnss/salary \\
6 & regime-service & Spring Boot & 8083 & ghcr.io/cnss/regime \\
7 & affiliation-service & Spring Boot & 8084 & ghcr.io/cnss/affiliation \\
8 & debit-service & Spring Boot & 8085 & ghcr.io/cnss/debit \\
9 & payment-service & Spring Boot & 8086 & ghcr.io/cnss/payment \\
10 & notification-service & Spring Boot + Mail & 8087 & ghcr.io/cnss/notification \\
11 & file-service & Spring Boot & 8088 & ghcr.io/cnss/file \\
12 & disponibilite-service & Spring Boot & 8091 & ghcr.io/cnss/disponibilite \\
13 & ai-extraction-service & Python FastAPI & 8090 & ghcr.io/cnss/ai-extraction \\
14 & frontend-cooperation & Angular 17 (NGINX) & 80 & ghcr.io/cnss/frontend-coop \\
15 & frontend-disponibilite & Angular 17 (NGINX) & 80 & ghcr.io/cnss/frontend-dispo \\
\hline
\end{longtable}
\end{center}

\newpage

% ============================================
\section{Outils d'Intégration Continue (CI/CD)}
% ============================================

\subsection{Comparatif des outils}

\begin{center}
\rowcolors{2}{light}{white}
\begin{tabular}{|l|p{4cm}|p{3.5cm}|p{3.5cm}|}
\hline
\rowcolor{primary}
\textcolor{white}{\textbf{Outil}} & \textcolor{white}{\textbf{Points forts}} & \textcolor{white}{\textbf{Points faibles}} & \textcolor{white}{\textbf{Idéal pour}} \\
\hline
\textbf{GitHub Actions} \badge{primary}{Recommandé} & Intégré à GitHub, YAML simple, marketplace riche & Minutes limitées (plan gratuit) & Projets GitHub \\
\textbf{Jenkins} & Très extensible, plugins nombreux, gratuit & Configuration complexe, UI vieillissante & Entreprises \\
\textbf{Azure DevOps} & Intégration Azure/MS parfaite & Vendor lock-in & Environnements MS \\
\hline
\end{tabular}
\end{center}

\subsection{Pourquoi GitHub Actions ?}

\begin{itemize}[leftmargin=*]
    \item \textbf{Intégré à GitHub} --- pas de serveur séparé à gérer
    \item \textbf{YAML simple} --- fichiers \texttt{.github/workflows/*.yml}
    \item \textbf{Marketplace riche} --- actions prêtes à l'emploi
    \item \textbf{Runners gratuits} --- 2000 minutes/mois (plan gratuit)
    \item \textbf{Self-hosted runners} --- possibilité d'utiliser ses propres serveurs
\end{itemize}

\newpage

% ============================================
\section{Pipeline GitHub Actions}
% ============================================

\subsection{Diagramme du pipeline}

\begin{center}
    \fbox{\includesvg[width=0.95\textwidth]{images/pipeline-cicd.svg}}
\end{center}

\subsection{Secrets GitHub à configurer}

\begin{center}
\rowcolors{2}{light}{white}
\begin{tabular}{|l|l|l|}
\hline
\rowcolor{primary}
\textcolor{white}{\textbf{Secret}} & \textcolor{white}{\textbf{Description}} & \textcolor{white}{\textbf{Exemple}} \\
\hline
\texttt{SONAR\_TOKEN} & Token SonarQube & \texttt{sqp\_xxxxx} \\
\texttt{SONAR\_HOST\_URL} & URL serveur SonarQube & \texttt{https://sonar.cnss.tn} \\
\texttt{KUBECONFIG\_STAGING} & Kubeconfig base64 (staging) & \texttt{base64 ~/.kube/config} \\
\texttt{KUBECONFIG\_PROD} & Kubeconfig base64 (production) & \texttt{base64 ~/.kube/config} \\
\texttt{SLACK\_WEBHOOK\_URL} & Webhook Slack & \texttt{https://hooks.slack.com/...} \\
\hline
\end{tabular}
\end{center}

\subsection{Extrait du workflow}

\begin{lstlisting}[language=bash,caption={.github/workflows/ci-cd.yml}]
name: CNSS CI/CD Pipeline

on:
  push:
    branches: [main, develop]
  pull_request:
    branches: [main]

env:
  REGISTRY: ghcr.io
  IMAGE_PREFIX: ${{ github.repository_owner }}/cnss

jobs:
  build-backend:
    runs-on: ubuntu-latest
    strategy:
      matrix:
        service: [auth-service, employer-service, salary-service]
    steps:
      - uses: actions/checkout@v4
      - uses: actions/setup-java@v4
        with:
          java-version: '17'
      - run: ./gradlew clean build test
\end{lstlisting}

\newpage

% ============================================
\section{Orchestration Kubernetes}
% ============================================

\subsection{Comparatif des solutions}

\begin{center}
\rowcolors{2}{light}{white}
\begin{tabular}{|l|p{4cm}|p{3.5cm}|l|}
\hline
\rowcolor{primary}
\textcolor{white}{\textbf{Solution}} & \textcolor{white}{\textbf{Avantages}} & \textcolor{white}{\textbf{Inconvénients}} & \textcolor{white}{\textbf{Coût}} \\
\hline
\textbf{Kubernetes (k3s)} \badge{primary}{Recommandé} & Léger, rapide, prod-ready & Moins de features que K8s complet & Gratuit \\
\textbf{Docker Swarm} & Simple, intégré à Docker & Moins de fonctionnalités & Gratuit \\
\textbf{Portainer} & UI intuitive, gère Docker/K8s & Limité pour gros clusters & Freemium \\
\textbf{OpenShift} & K8s + sécurité + CI/CD intégrés & Coûteux, lourd & Licence Red Hat \\
\hline
\end{tabular}
\end{center}

\subsection{Architecture du cluster}

\begin{center}
    \fbox{\includesvg[width=0.95\textwidth]{images/kubernetes-cluster.svg}}
\end{center}

\subsection{Installation k3s}

\begin{lstlisting}[language=bash,caption={Installation k3s sur les nodes}]
# Sur le master
curl -sfL https://get.k3s.io | sh -

# Recuperer le token
cat /var/lib/rancher/k3s/server/node-token

# Sur les workers
curl -sfL https://get.k3s.io | K3S_URL=https://master:6443 \
  K3S_TOKEN=<token> sh -
\end{lstlisting}

\newpage

% ============================================
\section{Monitoring et Logging}
% ============================================

\subsection{Stack de monitoring}

\begin{center}
    \fbox{\includesvg[width=0.95\textwidth]{images/monitoring-stack.svg}}
\end{center}

\subsection{Installation via Helm}

\begin{lstlisting}[language=bash,caption={Installation Prometheus + Grafana + Loki}]
# Ajouter les repos Helm
helm repo add prometheus-community \
  https://prometheus-community.github.io/helm-charts
helm repo add grafana https://grafana.github.io/helm-charts

# Installer Prometheus + Grafana
helm install prometheus prometheus-community/kube-prometheus-stack \
  --namespace monitoring --create-namespace

# Installer Loki pour les logs
helm install loki grafana/loki-stack --namespace monitoring
\end{lstlisting}

\subsection{Composants}

\begin{center}
\rowcolors{2}{light}{white}
\begin{tabular}{|l|l|l|}
\hline
\rowcolor{primary}
\textcolor{white}{\textbf{Composant}} & \textcolor{white}{\textbf{Fonction}} & \textcolor{white}{\textbf{Port}} \\
\hline
Prometheus & Collecte des métriques & 9090 \\
Alertmanager & Gestion des alertes & 9093 \\
Grafana & Dashboards & 3000 \\
Promtail & Collecte des logs & - \\
Loki & Stockage des logs & 3100 \\
\hline
\end{tabular}
\end{center}

\newpage

% ============================================
\section{Sécurité en Production}
% ============================================

\subsection{Checklist sécurité}

\begin{center}
\rowcolors{2}{light}{white}
\begin{tabular}{|c|l|p{6cm}|l|}
\hline
\rowcolor{primary}
\textcolor{white}{\textbf{\#}} & \textcolor{white}{\textbf{Élément}} & \textcolor{white}{\textbf{Action}} & \textcolor{white}{\textbf{Statut}} \\
\hline
1 & Secrets & Kubernetes Secrets ou HashiCorp Vault & \badge{warning}{Requis} \\
2 & Images Docker & Scanner avec Trivy ou Snyk & \badge{warning}{Requis} \\
3 & Registre Docker & ghcr.io avec authentification & \badge{success}{OK} \\
4 & HTTPS/TLS & cert-manager + Let's Encrypt & \badge{warning}{Requis} \\
5 & RBAC & Limiter permissions par namespace & \badge{warning}{Requis} \\
6 & Network Policies & Restreindre trafic entre pods & \badge{secondary}{Recommandé} \\
\hline
\end{tabular}
\end{center}

\subsection{Scanner les images avec Trivy}

\begin{lstlisting}[language=bash,caption={Scan de sécurité des images Docker}]
# Installer Trivy
curl -sfL https://raw.githubusercontent.com/aquasecurity/trivy/main/contrib/install.sh | sh -s -- -b /usr/local/bin

# Scanner une image
trivy image ghcr.io/cnss/auth-service:latest

# Scanner avec seuil de severite
trivy image --severity HIGH,CRITICAL ghcr.io/cnss/auth-service:latest
\end{lstlisting}

\newpage

% ============================================
\section{Solutions Cloud}
% ============================================

\subsection{Comparatif des options}

\begin{center}
\rowcolors{2}{light}{white}
\begin{tabular}{|l|l|l|p{4cm}|}
\hline
\rowcolor{primary}
\textcolor{white}{\textbf{Solution}} & \textcolor{white}{\textbf{Coût mensuel}} & \textcolor{white}{\textbf{Complexité}} & \textcolor{white}{\textbf{Idéal pour}} \\
\hline
\textbf{On-Premise (k3s)} \badge{primary}{Recommandé} & Coût serveurs & Moyen & Budget limité, contrôle total \\
AWS EKS & ~150-500\$/mois & Faible & Scalabilité, intégrations AWS \\
Azure AKS & ~100-400\$/mois & Faible & Environnement Microsoft \\
Google GKE & ~100-400\$/mois & Faible & Performance, ML/IA \\
OVHcloud Managed K8s & ~80-300\$/mois & Faible & Souveraineté données (EU) \\
\hline
\end{tabular}
\end{center}

\begin{infobox}
\textbf{Recommandation pour CNSS (Tunisie)}

\textbf{Option 1 :} Infrastructure locale avec 3 serveurs (16 Go RAM, 4 vCPU) + k3s + Portainer

\textbf{Option 2 :} Cloud hybride avec OVHcloud/Scaleway + Oracle DB on-premise + VPN
\end{infobox}

\newpage

% ============================================
\section{Résumé des outils recommandés}
% ============================================

\begin{center}
\rowcolors{2}{light}{white}
\begin{tabular}{|l|l|l|}
\hline
\rowcolor{primary}
\textcolor{white}{\textbf{Catégorie}} & \textcolor{white}{\textbf{Outil principal}} & \textcolor{white}{\textbf{Alternative}} \\
\hline
CI/CD & \badge{primary}{GitHub Actions} & Jenkins \\
Qualité code & \badge{secondary}{SonarQube} & --- \\
Registre Docker & \badge{success}{ghcr.io} & Docker Hub / Harbor \\
Orchestration & \badge{primary}{Kubernetes (k3s)} & Docker Swarm \\
Gestion K8s & \badge{secondary}{Portainer} & Rancher \\
Déploiement & \badge{success}{Helm + ArgoCD} & kubectl apply \\
Monitoring & \badge{warning}{Prometheus + Grafana} & Datadog (payant) \\
Logs & \badge{secondary}{Loki + Grafana} & ELK Stack \\
Sécurité images & \badge{primary}{Trivy} & Snyk \\
Certificats SSL & \badge{success}{cert-manager + Let's Encrypt} & --- \\
\hline
\end{tabular}
\end{center}

\vspace{1cm}

\begin{successbox}
\textbf{Stack recommandée pour CNSS}

\texttt{GitHub Actions + k3s + Portainer + Helm + Prometheus/Grafana + Loki}

Cette stack est \textbf{gratuite}, \textbf{robuste}, et adaptée à une administration publique.
\end{successbox}

\vfill

\begin{center}
\textcolor{gray}{\rule{0.5\textwidth}{0.4pt}}

\textbf{Guide DevOps --- Mise en Production CNSS}

Coopération Technique ATCT + Mise en Disponibilité Spéciale

Version 1.0 --- Mars 2026
\end{center}

\end{document}
